% Sinais Digitais e Analógicos

% Introdução - requisitos

% Amostragem
% 	O Teorema da Amostragem (Nyquist-Shannon)
% 	Aliasing (Subamostragem)
% 	Filtro Anti-Aliasing 

% Parâmetros fixos

% Método Ideal:
% - Descrição Analítica
% - Simulação de amostragem
% - Filtro de reconstrução
% - Simulação de reconstrução

% Método Natural:
% - Descrição Analítica
% - Simulação de amostragem
% - Filtro de reconstrução
% - Simulação de reconstrução

% Método Topo-Plano:
% - Descrição Analítica
% - Simulação de amostragem
% - Filtro de reconstrução
% - Simulação de reconstrução

% Cenários de Teste:
% - Casos de aliasing:
% 	- Variando f de corte e f de amostragem
% - Tau:
% 	- Aumentar/diminuir e explicar as consequências
% - Vetor de frequências:
% 	- Aumentar e diminuir o vetor para demostrar a sobreposições de sincs

% Resultados e Discussões:
% - Comparação entre métodos
% - Comparar em tempo e frequência
% 
% Scripts não precisam ser mostrados na apresentação, mas vão ser entregues e podem ser citados em questionamentos. GRAFICO DE MAGNITUDE E FASER PARA TOPO PLANO


\section{Introdução}

\begin{frame}{Contextualização}

    A empresa \textit{SigmaDSP Inc.} enfrenta o desafio de compreender minuciosamente o \textbf{teorema de amostragem} de sinais analógicos de tempo contínuo. Para isso sabe-se que três são as formas ou
métodos da amostragem de sinais analógicos de tempo contínuo:
    \begin{itemize}
        \item Ideal
        \item Natural 
        \item Instantânea de topo plano \textit{(flat-top)}.
    \end{itemize}
    Partindo desse contexto, foram solicitadas as simulações dos \textbf{métodos de amostragem} para uma senoide e os respectivos filtros de reconstrução. 
    
\end{frame}

\begin{frame}{Requisitos}

	Os requisitos para a resolução do problema eram:
	\begin{itemize}
		\item Simulação dos métodos de amostragem com uma senoide 
            \begin{itemize}
                \item \textit{Com uma segunda senoide sendo utilizada como impulso de frequência acima da máxima permitida pelo filtro.}
            \end{itemize}
        \item Frequência de amostragem ajustável 
            \begin{itemize}
                \item Para conseguir demonstrar o fenômeno do \textit{\textbf{aliasing}}
            \end{itemize}
        \item Filtragem de reconstrução
        \item Descrição e comparação analítica de todos os métodos e suas reconstruções. 
        
		
	\end{itemize}
\end{frame}